\section{Case definition}
\label{sec:case}

\subsection{The noses}

\begin{table}[H]
\centering
\caption{The two noses of report~2. Lengths in mm; $x$ is measured from the rounded tip.}
\label{tab:noses}
\small
\begin{tabular}{@{}lll@{}}
\toprule
 & Study A (this report) & Study B (to be run) \\
\midrule
Profile & LV-Haack & Von K\'arm\'an \\
Nominal curved length $L$, fineness $L/D$ & 561, 3.117 & 630, 3.5 \\
Coupler length $c$, total length & 89, 650 & 20, 650 \\
Tip radius $r_n$ & 4.5 & 4.5 \\
End of the curved part (shoulder), $x$ & 558.65 & 625.42 \\
End of the coupler, $x$ & 647.65 & 645.42 \\
\bottomrule
\end{tabular}
\end{table}

\noindent Both noses are solids of revolution in mm, with $x$ along the axis and the tip at $x = 0$. The
\SI{4.5}{mm} tip rounding shortens the curved part by \SI{2.35}{mm} (A) and \SI{4.58}{mm} (B). STEP files:
\file{CFD/geometry_650/nose_lvhaack_L561_c89_rn4.5mm.step} (A),
\file{nose_vonkarman_L630_c20_rn4.5mm.step} (B).

\subsection{Flight conditions}

\begin{table}[H]
\centering
\caption{Flight conditions of the three runs (altitude above sea level).}
\label{tab:points}
\small
\begin{tabular}{@{}rrrrl@{}}
\toprule
$M$ & altitude & $p_\infty$ & $T_\infty$ & type of flow \\
\midrule
0.95 & \SI{928}{m}  & \SI{90666}{Pa} & \SI{282.1}{K} & transonic: just below the speed of sound \\
1.20 & \SI{1478}{m} & \SI{84789}{Pa} & \SI{278.5}{K} & supersonic, close to $M = 1$ \\
1.47 & \SI{2189}{m} & \SI{77649}{Pa} & \SI{273.9}{K} & supersonic, top speed of the flight \\
\bottomrule
\end{tabular}
\end{table}

\noindent The inputs are $M$, $p_\infty$ and $T_\infty$, from the report-1 trajectory. On the report-2
trajectory of nose~A the same Mach numbers are reached within \SI{14}{m} of altitude and 0.2\,\% in
pressure; its top speed is $M = 1.457$ (\SI{2182}{m}), so the $M = 1.47$ run is the top-speed point
($+1\,\%$).

\paragraph{Why the speed matters.} Pressure waves travel upstream at the speed $a - V$ relative to the
rocket:
\begin{itemize}
\item $M < 1$: $a - V > 0$, the air ahead feels the nose coming. No bow shock; a small supersonic zone
forms at the shoulder and ends in a weak shock (transonic flow).
\item $M > 1$: $a - V < 0$, the air ahead is undisturbed until a bow shock just in front of the tip. The
high pressure behind it, acting on the front of the nose, is the wave drag.
\end{itemize}

\paragraph{What the runs give.} The drag on \zone{wall\_nose} (tip to shoulder), split by Fluent into
pressure and viscous parts:
\begin{equation}
C_D = \frac{F_x}{q_\infty A_{\rm ref}}, \qquad q_\infty = \tfrac12\rho_\infty V_\infty^2, \qquad
A_{\rm ref} = \frac{\pi D^2}{4}.
\end{equation}

\subsection{Symbols}

\begin{small}
\noindent
$a$ speed of sound; $A_{\rm ref}$ reference area $\pi D^2/4$; $C_D$ drag coefficient; $C_f$ friction
coefficient; $C_p$ pressure coefficient; CFL Courant number; $c_p$ specific heat at constant pressure; $D$
body diameter (\SI{180}{mm}); $F_x$ force along the rocket axis; $h_1$ height of the first cell at the
wall; $I$ turbulence intensity; $k$ turbulent kinetic energy; $M$ Mach number; $p$ pressure; $T_0$
stagnation (total) temperature; $q = \tfrac12\rho V^2$ dynamic pressure; $R$ gas constant of
air; $\mathrm{Re}$ Reynolds number; $r_n$ tip radius; $T$ temperature; $u_\tau$ friction velocity; $V$
flow speed; $y^+$ wall distance in viscous units; $\gamma$ ratio of specific heats; $\delta$
boundary-layer thickness; $\mu$ viscosity of the air; $\mu_t$ turbulent (eddy) viscosity; $\rho$
density; $\tau_w$ wall shear stress; $\omega$ turbulence frequency (specific dissipation rate).
Subscript $\infty$: the undisturbed air far from the rocket (free stream).
\end{small}
